% Copyright (c) 2025 Intel Corporation.  All rights reserved.
% SPDX-License-Identifier: Apache-2.0
%
% photopic-manual.tex -- User Manual for the Photopic Luminous Efficacy Optimizer

\documentclass[11pt]{article}

\usepackage[margin=1in]{geometry}
\usepackage{listings}
\usepackage{xcolor}
\usepackage{hyperref}
\usepackage{booktabs}
\usepackage{amsmath}
\usepackage{microtype}

\hypersetup{
  colorlinks=true,
  linkcolor=blue!60!black,
  urlcolor=blue!60!black,
  citecolor=blue!60!black,
}

\lstdefinestyle{shell}{
  language=bash,
  basicstyle=\ttfamily\small,
  backgroundcolor=\color{gray!10},
  frame=single,
  framerule=0.4pt,
  rulecolor=\color{gray!50},
  breaklines=true,
  columns=fullflexible,
  keepspaces=true,
  showstringspaces=false,
  xleftmargin=2em,
  framexleftmargin=1.5em,
  numbers=none,
  morekeywords={photopic,cm3,python3,cd,make},
}

\lstdefinestyle{scheme}{
  language=Lisp,
  basicstyle=\ttfamily\small,
  backgroundcolor=\color{blue!3},
  frame=single,
  framerule=0.4pt,
  rulecolor=\color{blue!30},
  breaklines=true,
  columns=fullflexible,
  keepspaces=true,
  showstringspaces=false,
  xleftmargin=2em,
  framexleftmargin=1.5em,
  numbers=none,
  keywordstyle=\color{blue!70!black}\bfseries,
  commentstyle=\color{green!50!black}\itshape,
  stringstyle=\color{red!60!black},
  morekeywords={define,lambda,let,let*,begin,if,cond,set!,and,or},
}

\lstdefinestyle{output}{
  basicstyle=\ttfamily\footnotesize,
  backgroundcolor=\color{yellow!5},
  frame=single,
  framerule=0.4pt,
  rulecolor=\color{yellow!50!black},
  breaklines=true,
  columns=fullflexible,
  keepspaces=true,
  showstringspaces=false,
  xleftmargin=2em,
  framexleftmargin=1.5em,
  numbers=none,
}

\newcommand{\photopic}{\texttt{photopic}}
\newcommand{\cmd}[1]{\texttt{#1}}
\newcommand{\file}[1]{\texttt{#1}}

\title{\textbf{Photopic} \\[0.3em]
  \large A Luminous Efficacy Optimizer for Spectral Power Distributions}
\author{User Manual}
\date{March 2026}

\begin{document}

\maketitle

\tableofcontents
\newpage

%======================================================================
\section{Introduction}
\label{sec:introduction}

\photopic{} is a constrained optimization tool that finds spectral power
distributions (SPDs) maximizing the \emph{luminous efficacy of radiation}
(LER) or the $R_9$ color rendering index, subject to constraints on CRI
(Color Rendering Index), CCT (Correlated Color Temperature), $D_{uv}$
(distance from the Planckian locus), and individual color rendering
components $R_i$.

The tool answers questions of the form: ``What is the theoretical maximum
luminous efficacy achievable at 2700\,K with CRI $\geq$ 80?'' and,
conversely, ``What is the best $R_9$ (saturated red rendering) achievable
at 2700\,K with CRI $\geq$ 80 and efficacy $\geq$ 350\,lm/W?''

\subsection{Background}

The luminous efficacy of radiation (LER) measures how efficiently a
spectrum produces visible light, in lumens per watt of radiated power.
The theoretical maximum is 683\,lm/W, achieved by monochromatic light at
555\,nm---which has a CRI of zero.  Real light sources must balance
efficacy against color quality: broader spectra render colors more
faithfully but waste energy in wavelength regions where the eye is less
sensitive.

\photopic{} explores this trade-off computationally, using the COBYLA
derivative-free constrained optimizer to search a parametric space of
spectral shapes.  Color rendering is evaluated via the CIE~13.3 test
color method with 14 test color samples (TCS).

\subsection{How it works}

\photopic{} starts from a blackbody spectrum at the target CCT and applies
a multiplicative modifier of the form $\exp(f(\lambda))$, where
$f(\lambda)$ is a piecewise-linear function defined by a parameter vector.
The optimizer adjusts these parameters to minimize (or maximize) the
objective while satisfying the constraints.

To avoid poor local minima in the high-dimensional space, the optimization
uses \emph{hierarchical subdivision}: it begins with 2 parameters (a
single linear slope across the visible spectrum), optimizes, then doubles
the resolution to 3, 5, 9, 17, 33, 65, and finally 129 parameters
(approximately 3\,nm spacing across 380--770\,nm).  At each level, the
COBYLA step size is halved.

%======================================================================
\section{Mathematical Foundations}
\label{sec:math}

This section summarizes the physics and colorimetry underlying the
optimization.  For a full treatment and computational results, see the
companion report~\cite{nystrom2026}.

\subsection{Planck's law and the blackbody reference}

The spectral radiance of a blackbody at temperature $T$ is
\begin{equation}
  B(\lambda, T) = \frac{2hc^2}{\lambda^5}
    \cdot \frac{1}{\exp\!\bigl(hc / \lambda k T\bigr) - 1}
  \label{eq:planck}
\end{equation}
where $h$ is Planck's constant, $c$ the speed of light, and $k$
Boltzmann's constant.  This serves as the starting point for the
optimization: the parametric spectrum is a multiplicative perturbation
of $B(\lambda, T)$.

\subsection{Luminous efficacy of radiation}

The luminous efficacy of radiation (LER) measures how effectively a
spectral power distribution $S(\lambda)$ produces the sensation of
brightness:
\begin{equation}
  \eta = K_m \,
    \frac{\displaystyle\int_{\lambda_0}^{\lambda_1}
      S(\lambda)\, V(\lambda)\, d\lambda}
    {\displaystyle\int_{\lambda_0}^{\lambda_1}
      S(\lambda)\, d\lambda}
  \label{eq:ler}
\end{equation}
where $V(\lambda)$ is the CIE photopic luminous efficiency function
(peaking at unity at 555\,nm), $K_m = 683$\,lm/W is the maximum
luminous efficacy, and $[\lambda_0, \lambda_1] = [380, 770]$\,nm
spans the visible spectrum.  For monochromatic light at 555\,nm,
$\eta = 683$\,lm/W; for a broadband source, $\eta$ is necessarily
lower.  \emph{This is the primary quantity that \photopic{} maximizes.}

\subsection{Tristimulus values and chromaticity}

The CIE 1931 tristimulus values of a spectrum $S(\lambda)$ are
\begin{equation}
  X = \int_{\lambda_0}^{\lambda_1} S(\lambda)\, \bar{x}(\lambda)\, d\lambda, \quad
  Y = \int_{\lambda_0}^{\lambda_1} S(\lambda)\, \bar{y}(\lambda)\, d\lambda, \quad
  Z = \int_{\lambda_0}^{\lambda_1} S(\lambda)\, \bar{z}(\lambda)\, d\lambda
  \label{eq:tristimulus}
\end{equation}
where $\bar{x}$, $\bar{y}$, $\bar{z}$ are the CIE 1931 $2^\circ$
standard observer color matching functions, tabulated at 1\,nm
resolution in \file{CieXyz.i3}.  The chromaticity coordinates are
\begin{equation}
  x = \frac{X}{X+Y+Z}, \qquad y = \frac{Y}{X+Y+Z}
\end{equation}
and the CIE 1960 UCS coordinates used for CCT determination are
\begin{equation}
  u = \frac{4X}{X + 15Y + 3Z}, \qquad
  v = \frac{6Y}{X + 15Y + 3Z}.
  \label{eq:ucs}
\end{equation}

\subsection{Correlated color temperature}

The correlated color temperature (CCT) of a source is the temperature
$T$ of the Planckian radiator whose chromaticity is nearest in
$(u,v)$ space:
\begin{equation}
  \text{CCT} = \operatorname*{arg\,min}_{T} \;
    \bigl\lVert (u_{\text{source}}, v_{\text{source}})
    - (u_{\text{Planck}}(T), v_{\text{Planck}}(T)) \bigr\rVert
  \label{eq:cct}
\end{equation}
The signed distance from the Planckian locus is $D_{uv}$, constrained
to $D_{uv} \leq 0.012$ to ensure the source chromaticity lies
acceptably close to the blackbody locus.  \photopic{} determines the
CCT by Brent's method search over a precomputed temperature-to-$uv$
lookup table (\file{TempUv.m3}).

\subsection{Color Rendering Index}

The CIE General Color Rendering Index $R_a$ quantifies how faithfully
a test illuminant renders colors compared to a reference illuminant
(a blackbody at the same CCT).  For each of 14~test color samples
(TCS) with reflectance $\rho_i(\lambda)$, the reflected spectrum
under the test illuminant is
\begin{equation}
  S_i^{\text{test}}(\lambda) =
    \hat{S}(\lambda) \cdot \rho_i(\lambda)
  \label{eq:reflected}
\end{equation}
where $\hat{S}$ is the Y-normalized test spectrum ($Y = 100$).
The reflected spectrum is converted to CIE 1964 $U^*V^*W^*$
coordinates with von~Kries chromatic adaptation:
\begin{align}
  W^* &= 25\, Y^{1/3} - 17 \\
  U^* &= 13\, W^* \cdot (u' - u_0) \\
  V^* &= 13\, W^* \cdot (v' - v_0)
  \label{eq:uvw}
\end{align}
where $(u_0, v_0)$ is the adapted white point.  The color difference
between test and reference renderings is
\begin{equation}
  \Delta E_i = \sqrt{
    (\Delta U^*_i)^2 + (\Delta V^*_i)^2 + (\Delta W^*_i)^2
  }
\end{equation}
and the special color rendering index for sample $i$ is
\begin{equation}
  R_i = 100 - 4.6 \, \Delta E_i.
  \label{eq:ri}
\end{equation}
The general index averages the first eight (unsaturated pastel) samples:
\begin{equation}
  R_a = \frac{1}{8} \sum_{i=1}^{8} R_i.
  \label{eq:cri}
\end{equation}
Notably, $R_9$ (saturated red) through $R_{14}$ are \emph{not}
included in $R_a$.  This is the well-known ``CRI loophole'': a
spectrum can score $R_a \geq 80$ while rendering saturated reds
poorly ($R_9 < 0$), which is precisely the scenario explored by the
inverse problem mode of \photopic{}.

\subsection{The optimization problem}

In the forward mode (\cmd{-run}), \photopic{} solves:
\begin{equation}
\fboxsep=8pt
\boxed{\;
  \max_{\mathbf{p}} \; \eta\bigl(S(\lambda; \mathbf{p})\bigr)
  \quad \text{subject to} \quad
  \begin{cases}
    R_a \geq R_a^{\min} \\[3pt]
    \min_{i \in 1..8} R_i \geq R_a^{\min} - 10 \\[3pt]
    R_9 \geq R_9^{\min} \\[3pt]
    |\text{CCT} - T_{\text{target}}| \leq 50\,\text{K} \\[3pt]
    D_{uv} \leq 0.012
  \end{cases}
\;}
  \label{eq:forward}
\end{equation}

\medskip\noindent
In the inverse mode (\cmd{-run-maxr9}):
\begin{equation}
\fboxsep=8pt
\boxed{\;
  \max_{\mathbf{p}} \; R_9\bigl(S(\lambda; \mathbf{p})\bigr)
  \quad \text{subject to} \quad
  \begin{cases}
    R_a \geq R_a^{\min} \\[3pt]
    \min_{i \in 1..8} R_i \geq R_a^{\min} - 10 \\[3pt]
    \eta \geq \eta_{\min} \\[3pt]
    |\text{CCT} - T_{\text{target}}| \leq 50\,\text{K} \\[3pt]
    D_{uv} \leq 0.012
  \end{cases}
\;}
  \label{eq:inverse}
\end{equation}
Both are solved by COBYLA, Powell's derivative-free constrained
optimizer~\cite{powell1994}, which constructs linear approximations
to the objective and constraint functions within a shrinking
trust region.

\subsection{Parametric spectrum}

The spectrum is parameterized as
\begin{equation}
  S(\lambda; \mathbf{p}) =
    B(\lambda, T) \cdot
    \exp\!\biggl(\,
      \sum_{j} p_j \, \phi_j(\lambda)
    \biggr)
  \label{eq:parametric}
\end{equation}
where $\phi_j(\lambda)$ are hat functions (piecewise-linear basis
functions) centered at uniformly spaced nodes
$\lambda_j = \lambda_0 + j \cdot \Delta\lambda$ with
$\Delta\lambda = (\lambda_1 - \lambda_0) / (n-1)$.  The exponent
$f(\lambda) = \sum_j p_j \phi_j(\lambda)$ is the linear interpolant
of the parameter values $p_j$ at the nodes.  At $\mathbf{p} = \mathbf{0}$,
$S = B$: the unperturbed blackbody.  The exponential ensures
$S(\lambda) > 0$ for all $\mathbf{p}$.

The hierarchical subdivision doubles the number of nodes at each stage,
interpolating the coarse solution onto the fine grid:
\begin{align}
  n &\in \{2,\; 3,\; 5,\; 9,\; 17,\; 33,\; 65,\; 129\} \\[4pt]
  \Delta\lambda &\in \{390,\; 195,\; 97.5,\; \ldots,\; 3.05\} \;\text{nm}.
\end{align}

\begin{thebibliography}{9}

\bibitem{nystrom2026}
M.~Nystr\"om,
\emph{Theoretical Limits of Luminous Efficacy under Color Rendering
Constraints: A Computational Study},
Technical Report, 2026.
Available at \url{https://github.com/IntelLabs/async-toolkit}, branch
\texttt{claude\_edits}.

\bibitem{powell1994}
M.\,J.\,D.~Powell,
``A direct search optimization method that models the objective and
constraint functions by linear interpolation,''
in \emph{Advances in Optimization and Numerical Analysis},
S.~Gomez and J.-P.~Hennart, eds., pp.~51--67, Kluwer, 1994.

\bibitem{thornton1971}
W.\,A.~Thornton,
``Luminosity and color-rendering capability of white light,''
\emph{J.\ Opt.\ Soc.\ Am.}, vol.~61, no.~9, pp.~1155--1163, 1971.

\bibitem{walter1978}
W.~Walter,
``An optimization approach to the computation of luminous efficacy
bounds for sources with color rendering constraints,''
\emph{Color Res.\ Appl.}, vol.~3, no.~2, pp.~83--90, 1978.

\end{thebibliography}

%======================================================================
\section{Building}
\label{sec:building}

\subsection{Prerequisites}

\begin{itemize}
  \item A working CM3 (Critical Mass Modula-3) compiler, installed and
    in your \cmd{PATH} (or referenced by absolute path).
  \item The \cmd{m3utils} tree built---\photopic{} depends on several
    libraries within the tree including \cmd{mscheme},
    \cmd{modula3scheme}, \cmd{parseparams}, \cmd{integrate},
    \cmd{cobyla}, \cmd{minimize}, \cmd{mpfr}, and \cmd{cit\_util}.
  \item The \cmd{sstubgen} tool must also be available, since the build
    generates Scheme stubs for the Modula-3 interfaces.
\end{itemize}

\subsection{Build procedure}

Build the entire \cmd{m3utils} tree first (if not already done), then
build \photopic{}:

\begin{lstlisting}[style=shell]
# Build the full m3utils tree (from the m3utils root)
cd m3utils
export CM3=yes
gmake

# Or build just photopic (after dependencies are built)
cd photopic/src
cm3 -override
\end{lstlisting}

The binary is placed in a target-specific directory alongside the source,
e.g., \file{photopic/\allowbreak{}ARM64\_DARWIN/\allowbreak{}photopic}
on Apple Silicon macOS or
\file{photopic/\allowbreak{}AMD64\_LINUX/\allowbreak{}photopic}
on x86-64 Linux.

If \cmd{sstubgen} is not in your default \cmd{PATH}, prefix the build
command:

\begin{lstlisting}[style=shell]
env PATH="/path/to/cm3/install/bin:${PATH}" cm3 -override
\end{lstlisting}

%======================================================================
\section{Command-Line Interface}
\label{sec:cli}

\photopic{} is a Modula-3 program that embeds MScheme, a Scheme
interpreter.  The main optimization logic lives in
\file{photopic.scm}, which is loaded at startup via the \cmd{-scmfile}
flag.  All command-line invocations require the Scheme mode flags:

\begin{lstlisting}[style=shell]
photopic <mode> <args...> -scm -scmfile photopic.scm
\end{lstlisting}

The \cmd{-scm} flag activates the embedded Scheme interpreter.
The \cmd{-scmfile} flag specifies a Scheme source file to load (the path
to \file{photopic.scm} relative to the working directory or as an
absolute path).

\subsection{\texttt{-run}: Maximize efficacy}

\begin{lstlisting}[style=shell]
photopic -run <cct> <min-cri> <min-r9> -scm -scmfile photopic.scm
\end{lstlisting}

Maximizes luminous efficacy of radiation (LER, in lm/W) subject to:
\begin{itemize}
  \item CRI (Ra) $\geq$ \cmd{<min-cri>}
  \item $R_9$ $\geq$ \cmd{<min-r9>}
  \item Worst individual $R_i$ ($i \in 1..8$) $\geq$ \cmd{<min-cri>}$-10$
  \item CCT within $\pm50$\,K of \cmd{<cct>}
  \item $D_{uv} \leq 0.012$
\end{itemize}

The optimizer runs 8 hierarchical levels (2 to 129 parameters).
Set \cmd{<min-r9>} to \cmd{-100} to leave $R_9$ effectively unconstrained.

\subsection{\texttt{-run-maxr9}: Maximize \texorpdfstring{$R_9$}{R9} (inverse problem)}

\begin{lstlisting}[style=shell]
photopic -run-maxr9 <cct> <min-cri> <min-efficacy> \
         -scm -scmfile photopic.scm
\end{lstlisting}

Maximizes the $R_9$ color rendering index subject to:
\begin{itemize}
  \item CRI (Ra) $\geq$ \cmd{<min-cri>}
  \item Worst individual $R_i$ $\geq$ \cmd{<min-cri>}$-10$
  \item Efficacy $\geq$ \cmd{<min-efficacy>} lm/W
  \item CCT within $\pm50$\,K of \cmd{<cct>}
  \item $D_{uv} \leq 0.012$
\end{itemize}

This is the inverse of \cmd{-run}: instead of finding how efficient a
spectrum can be, it finds how well saturated reds can be rendered at a
given efficiency floor.

\subsection{\texttt{-run-drop}: Multi-run with TCS axis dropping}

\begin{lstlisting}[style=shell]
photopic -run-drop <cct> <min-cri> -scm -scmfile photopic.scm
\end{lstlisting}

Performs multiple rounds of efficacy maximization.  After each round, the
test color sample axis with the worst $R_i$ is removed from the CRI
average, and the optimizer re-runs against the remaining axes.  This
explores how much additional efficacy can be gained by sacrificing
individual color rendering dimensions.

%======================================================================
\section{Worked Examples}
\label{sec:examples}

All examples assume you are in the \file{report-runs/} directory with
the \photopic{} binary at \file{../ARM64\_DARWIN/\allowbreak{}photopic}
(adjust the path for your platform) and \file{photopic.scm} at
\file{../src/\allowbreak{}photopic.scm}.

For brevity, we define:
\begin{lstlisting}[style=shell]
PHOTOPIC=../ARM64_DARWIN/photopic
SCMFILE=../src/photopic.scm
\end{lstlisting}

\subsection{Example 1: Basic run at CRI $\geq$ 80}

\begin{lstlisting}[style=shell]
$PHOTOPIC -run 2700 80 -100 -scm -scmfile $SCMFILE
\end{lstlisting}

This maximizes efficacy for a 2700\,K source with CRI $\geq$ 80 and no
$R_9$ constraint.  Expect the optimizer to produce a spectrum with
efficacy around 400--420\,lm/W, achieved by concentrating power in
spectral regions where the eye is most sensitive while maintaining
adequate color rendering.  Output files are written to the current
directory (see Section~\ref{sec:output}).

\subsection{Example 2: High CRI at CRI $\geq$ 95}

\begin{lstlisting}[style=shell]
$PHOTOPIC -run 2700 95 -100 -scm -scmfile $SCMFILE
\end{lstlisting}

With a tighter CRI constraint, the optimizer has less freedom to reshape
the spectrum.  Expect efficacy around 350--370\,lm/W---a significant
penalty compared to CRI $\geq$ 80, because the spectrum must remain closer
to a blackbody shape to render all 8 primary test colors faithfully.

\subsection{Example 3: Moderate CRI with \texorpdfstring{$R_9$}{R9} constraint}

\begin{lstlisting}[style=shell]
$PHOTOPIC -run 2700 82 50 -scm -scmfile $SCMFILE
\end{lstlisting}

This adds an $R_9 \geq 50$ constraint, requiring decent saturated red
rendering.  Without this constraint, the optimizer often sacrifices deep
red content (which the eye is insensitive to) for higher efficacy.  The
$R_9$ floor forces retention of spectral power near 630\,nm, reducing
achievable efficacy.

\subsection{Example 4: Inverse problem---maximize \texorpdfstring{$R_9$}{R9} at efficacy \texorpdfstring{$\geq$}{>=} 350}

\begin{lstlisting}[style=shell]
$PHOTOPIC -run-maxr9 2700 80 350 -scm -scmfile $SCMFILE
\end{lstlisting}

Asks: given a 2700\,K source that must achieve CRI $\geq$ 80 and
efficacy $\geq$ 350\,lm/W, what is the best $R_9$ possible?  This
traces the Pareto frontier between efficacy and red rendering quality.

\subsection{Example 5: Running the full CRI sweep}

The script \file{run-optimizations.sh} automates running at multiple CRI
levels:

\begin{lstlisting}[style=shell]
./run-optimizations.sh            # all levels: 60 70 80 82 85 90 95 98
./run-optimizations.sh 90         # only CRI >= 90
./run-optimizations.sh 80 85 90   # specific levels
\end{lstlisting}

\subsection{Example 6: Sweeping the efficacy--\texorpdfstring{$R_9$}{R9} frontier}

The script \file{run-maxr9-sweep.sh} sweeps efficacy from 300 to
430\,lm/W in steps of 10, maximizing $R_9$ at each level:

\begin{lstlisting}[style=shell]
./run-maxr9-sweep.sh
\end{lstlisting}

\subsection{Example 7: Plotting results}

After running optimizations, generate publication-quality figures:

\begin{lstlisting}[style=shell]
python3 plot_spectra.py
\end{lstlisting}

This produces PDF and PNG files for:
\begin{itemize}
  \item Blackbody vs.\ optimized spectrum comparison
  \item Multi-panel CRI comparison
  \item Convergence across resolution levels
  \item Overlay of all CRI-constrained spectra
  \item Efficacy vs.\ CRI trade-off curve
  \item $R_9$ vs.\ efficacy scatter plot
\end{itemize}

%======================================================================
\section{Output Files}
\label{sec:output}

All output is written to the current working directory.

\subsection{File naming convention}

Output files follow the pattern:
\begin{center}
\cmd{<cct>\_CRI<min-cri>\_<suffix>\_<dims>.<ext>}
\end{center}
where:
\begin{itemize}
  \item \cmd{<cct>} is the target CCT (e.g., \cmd{2700})
  \item \cmd{<min-cri>} is the minimum CRI constraint (e.g., \cmd{82})
  \item \cmd{<suffix>} encodes run parameters (\cmd{R9=-100} for
    unconstrained $R_9$; \cmd{\_maxR9\_eff350} for the inverse problem)
  \item \cmd{<dims>} is the parameter count at the subdivision level
    (2, 3, 5, 9, 17, 33, 65, or 129)
\end{itemize}

One file is produced \emph{at each subdivision level}, so a complete run
at 8 levels generates 8 \file{.res} and 8 \file{.dat} files.

\subsection{\texttt{.res} files: Optimization results}

Each \file{.res} file is an S-expression on two lines.

\paragraph{Line 1: Computed specifications.}  A flat list with the
following fields:
\begin{center}
\begin{tabular}{@{}rll@{}}
\toprule
Position & Field & Description \\
\midrule
1  & CRI (Ra)    & Average of $R_1 \ldots R_8$ \\
2  & Worst $R_i$ & Minimum of $R_1 \ldots R_8$ \\
3  & Efficacy    & Luminous efficacy of radiation (lm/W) \\
4  & CCT         & Correlated color temperature (K) \\
5  & $D_{uv}$    & Distance from Planckian locus \\
6--19  & $R_1 \ldots R_{14}$ & Individual color rendering indices \\
20 & CRI-14      & Average of all 14 $R_i$ values \\
21 & Worst-14    & Minimum of all 14 $R_i$ values \\
22 & Sel-CRI     & CRI over selected TCS subset \\
23 & Sel-worst   & Worst $R_i$ over selected subset \\
24 & Worst axis  & Index of the worst TCS axis \\
\bottomrule
\end{tabular}
\end{center}

Example (from \file{2700\_CRI82\_R9=-100\_129.res}):

\begin{lstlisting}[style=output]
(83.22 72.01 408.61 (2650.00 0.012) (93.16 87.51
 72.01 92.19 80.80 75.68 92.36 72.02 28.67 57.20
 83.54 16.30 90.87 80.46) 73.06 16.30 73.06 16.30 12)
\end{lstlisting}

This result achieves 408.6\,lm/W at CRI\,=\,83.2 with CCT\,=\,2650\,K
and $R_9 = 28.7$.

\paragraph{Line 2: Parameter vector.}  The optimized parameter values
for the piecewise-linear log-modifier.  These can be used to reconstruct
the spectrum or as a warm-start for further optimization.

\subsection{\texttt{.dat} files: Spectral data}

Two types of \file{.dat} files are produced:

\begin{itemize}
  \item \file{w\_<name>\_<dims>.dat} --- the Y-normalized optimized
    spectrum
  \item \file{base\_<cct>.dat} --- the unmodified blackbody reference
    at temperature \cmd{<cct>}
\end{itemize}

Both are two-column text files (wavelength in meters, spectral power
density), with 391 data points at 1\,nm resolution spanning
380--770\,nm:

\begin{lstlisting}[style=output]
3.8e-7 5.366289953705953e6
3.8195e-7 5.7351389943960775e6
3.839e-7 6.1250440250444e6
...
\end{lstlisting}

To convert wavelengths to nanometers, multiply the first column by
$10^9$.

%======================================================================
\section{Architecture}
\label{sec:architecture}

\subsection{Parametric spectrum model}

The optimized spectrum is defined as:
\[
  S(\lambda) = B(\lambda, T) \cdot \exp\!\bigl(f(\lambda)\bigr)
\]
where $B(\lambda, T)$ is the Planck blackbody spectrum at the target CCT,
and $f(\lambda)$ is a piecewise-linear function defined by a parameter
vector $\mathbf{p} = (p_0, p_1, \ldots, p_{n-1})$.

The parameter nodes are uniformly spaced across 380--770\,nm.  Between
nodes, $f$ is linearly interpolated.  At $\mathbf{p} = \mathbf{0}$,
the spectrum equals the unmodified blackbody.  The exponential ensures
the spectrum remains strictly positive for all parameter values.

This representation is implemented in \file{ParametricSpectrum.m3}.
The function \cmd{GetFunc} returns an \cmd{LRFunction.T} that
evaluates $S(\lambda)$ given a parameter vector; \cmd{EvalFunction}
performs the interpolation lookup and exponential multiplication.

\subsection{Hierarchical subdivision}

Optimization begins at 2 parameters and proceeds through the sequence:
\[
  2 \to 3 \to 5 \to 9 \to 17 \to 33 \to 65 \to 129
\]
At each level, the number of parameters follows $n_{k+1} = 2n_k - 1$
(inserting a midpoint between every adjacent pair).  The
\cmd{Subdivide} procedure interpolates the old parameter vector into
the new, higher-resolution grid.  The COBYLA initial step size
(\cmd{rhobeg}) is halved at each level.

This hierarchical approach is essential: starting directly at 129
parameters would produce poor results because COBYLA would get trapped
in a local minimum near the blackbody starting point.  The
low-dimensional stages identify the dominant spectral features (which
wavelength regions to emphasize or suppress), and the subsequent stages
refine the details.

At 129 parameters, the node spacing is $(770 - 380) / 128 \approx 3.05$\,nm,
providing spectral features as narrow as approximately 6\,nm (two nodes
per feature).

\subsection{COBYLA optimizer}

COBYLA (Constrained Optimization BY Linear Approximation) is a
derivative-free optimizer by M.\,J.\,D.\ Powell.  It is well suited to
this problem because:
\begin{itemize}
  \item The CRI computation involves color adaptation and nonlinear
    coordinate transforms that are not easily differentiable.
  \item The constraint functions are nonlinear and involve search
    (CCT is found by minimizing distance in MacAdam $uv$ space).
  \item The number of function evaluations per optimization step scales
    linearly with the number of parameters.
\end{itemize}

The Modula-3 binding (\cmd{COBYLA\_M3.Minimize}) accepts the state
vector, number of constraints, an objective/constraint evaluation
function, \cmd{rhobeg} (initial trust-region radius), \cmd{rhoend}
(convergence tolerance), and a maximum iteration count.

\subsection{CRI calculation}

Color rendering is evaluated following CIE~13.3:
\begin{enumerate}
  \item Normalize the test spectrum so $Y = 100$.
  \item Determine the CCT by minimizing distance in MacAdam 1960 $uv$
    space to the Planckian locus (Brent's method via \cmd{search-T}).
  \item Generate the reference illuminant: a blackbody at the CCT.
  \item For each of 14 test color samples (TCS), compute the reflected
    spectrum under both the test and reference illuminants.
  \item Convert to CIE 1964 $U^*V^*W^*$ coordinates with von Kries
    chromatic adaptation.
  \item $R_i = 100 - 4.6 \, \Delta E_i$ where $\Delta E_i$ is the
    Euclidean distance between adapted and reference $U^*V^*W^*$ values.
  \item $\text{CRI (Ra)} = \frac{1}{8}\sum_{i=1}^{8} R_i$.
\end{enumerate}

The test color sample reflectance spectra are tabulated at 1\,nm
resolution in \file{Tcs.i3} (14 samples, covering pastels $R_1$--$R_8$
and saturated colors $R_9$--$R_{14}$).  CIE $\bar{x}$, $\bar{y}$,
$\bar{z}$ color matching functions are tabulated in \file{CieXyz.i3};
the photopic luminous efficiency function (683\,lm/W peak) is in
\file{CieSpectrum.i3}.

\subsection{Constraint convention}

In the objective/constraint function returned to COBYLA, the first
element of the returned vector is the quantity to \emph{minimize} (e.g.,
negative efficacy to maximize efficacy, or negative $R_9$ to maximize
$R_9$).  Subsequent elements are constraint values where
\textbf{positive means satisfied}:
\[
  c_i \geq 0 \quad \Longleftrightarrow \quad \text{constraint $i$ is met}
\]

For example, the CRI constraint returns $\text{CRI} - \text{min\_cri}$,
and the CCT constraint returns two values: $\text{CCT} - (\text{target} - 50)$
and $(\text{target} + 50) - \text{CCT}$.

%======================================================================
\section{Interactive Use}
\label{sec:interactive}

When invoked with \cmd{-scm -scmfile photopic.scm} but \emph{without}
a \cmd{-run}, \cmd{-run-maxr9}, or \cmd{-run-drop} flag, \photopic{}
drops into an interactive MScheme read-eval-print loop (REPL).
This is useful for exploration, debugging, and prototyping new objective
functions.

\begin{lstlisting}[style=shell]
$PHOTOPIC -scm -scmfile ../src/photopic.scm
\end{lstlisting}

\subsection{Useful interactive commands}

\paragraph{Setting up a problem.}

\begin{lstlisting}[style=scheme]
;; Create a parametric spectrum problem at 3000K with 5 parameters
(setup-problem!
  (ParametricSpectrum.MakeBlackbodyInWavelength 3000 l0 l1)
  5)
\end{lstlisting}

\paragraph{Inspecting the current spectrum.}

\begin{lstlisting}[style=scheme]
;; Compute full specs of the current spectrum
(calc-specs w)
;; => (cri-ra worst-ri efficacy (cct duv) (R1..R14) ...)

;; Get just the efficacy
(total-lumens-per-watt w)
\end{lstlisting}

\paragraph{Running one optimization pass.}

\begin{lstlisting}[style=scheme]
;; Run COBYLA with rhobeg=2 using the efficacy target
(run 2 specs->target)

;; Subdivide and run again at finer resolution
(subdivide-problem!)
(run 1 specs->target)
\end{lstlisting}

\paragraph{Evaluating a specific spectrum.}

\begin{lstlisting}[style=scheme]
;; CRI of a 4000K blackbody
(calc-cri (make-Bl 4000))

;; Efficacy of a fluorescent illuminant
(total-lumens-per-watt (FL 4))
\end{lstlisting}

\paragraph{Plotting.}

\begin{lstlisting}[style=scheme]
;; Write the current optimized spectrum to a gnuplot-compatible file
(plot (normalize-spectrum w) l0 l1 "my_spectrum.dat" 391)

;; Write the 14 TCS reflectance spectra
(plot-the-samples)
\end{lstlisting}

\paragraph{Modifying optimization parameters.}

\begin{lstlisting}[style=scheme]
;; Change target CCT
(set! *target-cct* 4000)

;; Change minimum CRI
(set! *min-cri-ra* 90)

;; Change minimum R9
(set! *min-r9* 50)

;; Change minimum efficacy (for -run-maxr9 mode)
(set! *min-efficacy* 350)
\end{lstlisting}

\paragraph{Reloading after edits.}

\begin{lstlisting}[style=scheme]
;; Reload photopic.scm after editing it
(reload)
\end{lstlisting}

%======================================================================
\section{Customization}
\label{sec:customization}

\subsection{Writing new objective functions}

The optimizer calls a \emph{spec-evaluator} function that takes the
output of \cmd{calc-specs} and returns a vector for COBYLA.
The existing evaluators follow this pattern:

\begin{lstlisting}[style=scheme]
(define (specs->my-target specs)
  (let ((lpm    (caddr specs))        ;; efficacy (lm/W)
        (r9     (nth (car (cddddr specs)) 8))  ;; R9
        (cri    (car specs))          ;; CRI (Ra)
        (crmin  (cadr specs))         ;; worst Ri (i in 1..8)
        (cct    (caar (cdddr specs))) ;; CCT (K)
        (Duv    (cadar (cdddr specs)));; Duv
        )
    (list <objective>        ;; minimize this (negate to maximize)
          <constraint-1>     ;; positive = satisfied
          <constraint-2>
          ...
          )))
\end{lstlisting}

The return value is a list whose first element is the objective to
minimize and whose remaining elements are constraint values (positive
means feasible).  Update \cmd{*num-constraints*} if you change the
number of constraints from the default~6.

\subsection{What \texorpdfstring{\texttt{calc-specs}}{calc-specs} returns}

The \cmd{calc-specs} function evaluates a spectrum and returns a list
with the following structure:

\begin{lstlisting}[style=scheme]
(cri-ra              ;; (car specs)      -- CRI average R1..R8
 worst-ri            ;; (cadr specs)     -- min of R1..R8
 efficacy            ;; (caddr specs)    -- lm/W
 ((cct duv)          ;; (cadddr specs)   -- CCT and Duv
  (R1 R2 ... R14))   ;; individual Ri values
 cri-14              ;; average of all 14 Ri
 worst-14            ;; min of all 14 Ri
 sel-cri             ;; CRI over selected subset
 sel-worst           ;; worst over selected subset
 worst-axis)         ;; index of worst TCS axis
\end{lstlisting}

\subsection{Accessing individual fields}

\begin{lstlisting}[style=scheme]
;; Given: (define s (calc-specs w))
(car s)                           ;; CRI (Ra)
(cadr s)                          ;; worst Ri
(caddr s)                         ;; efficacy
(caar (cdddr s))                  ;; CCT
(cadar (cdddr s))                 ;; Duv
(nth (car (cddddr s)) 8)          ;; R9 (9th of 14, 0-indexed)
(nth (car (cddddr s)) 0)          ;; R1
\end{lstlisting}

\subsection{Example: maximize efficacy subject to \texorpdfstring{$R_9 \geq 50$}{R9 >= 50}}

\begin{lstlisting}[style=scheme]
(define (specs->eff-r9floor specs)
  (let ((lpm  (caddr specs))
        (r9   (nth (car (cddddr specs)) 8))
        (cri  (car specs))
        (cct  (caar (cdddr specs)))
        (Duv  (cadar (cdddr specs))))
    (list (- lpm)                        ;; maximize efficacy
          (- cri 80)                     ;; CRI >= 80
          (- r9 50)                      ;; R9  >= 50
          (- cct (- *target-cct* 50))    ;; CCT in range
          (- (+ *target-cct* 50) cct)
          (* 1000 (- 0.012 Duv)))))      ;; Duv limit

;; Run it:
(set! *num-constraints* 5)
(run-example-iters! 2700 80 50 7 specs->eff-r9floor "_custom")
\end{lstlisting}

%======================================================================
\section{Reference}
\label{sec:reference}

\subsection{Global variables}

\begin{center}
\begin{tabular}{@{}lll@{}}
\toprule
Variable & Default & Purpose \\
\midrule
\cmd{*target-cct*}      & 2700  & Target correlated color temperature (K) \\
\cmd{*min-cri-ra*}      & 82    & Minimum CRI (Ra) constraint \\
\cmd{*min-r9*}          & $-100$ & Minimum $R_9$ constraint \\
\cmd{*min-efficacy*}    & 300   & Minimum efficacy (lm/W, for \cmd{-run-maxr9}) \\
\cmd{*num-constraints*} & 6     & Number of COBYLA constraints \\
\cmd{l0}                & $380 \times 10^{-9}$ & Blue limit of visible spectrum (m) \\
\cmd{l1}                & $770 \times 10^{-9}$ & Red limit of visible spectrum (m) \\
\bottomrule
\end{tabular}
\end{center}

\subsection{Key source files}

\begin{center}
\begin{tabular}{@{}ll@{}}
\toprule
File & Contents \\
\midrule
\file{src/photopic.scm}          & Main optimization logic (Scheme) \\
\file{src/Main.m3}               & Program entry point; Scheme REPL setup \\
\file{src/ParametricSpectrum.i3}  & Parametric spectrum interface \\
\file{src/ParametricSpectrum.m3}  & Parametric spectrum implementation \\
\file{src/CieSpectrum.i3}         & Photopic luminous efficiency data \\
\file{src/CieXyz.i3}              & CIE $\bar{x}\bar{y}\bar{z}$ color matching functions \\
\file{src/Tcs.i3}                 & Test color sample reflectance spectra (1--14) \\
\file{src/TempUv.m3}              & Temperature-to-$uv$ lookup table \\
\file{src/Blackbody.m3}           & Planck blackbody radiance \\
\file{src/m3makefile}             & Build dependencies and Scheme stub generation \\
\bottomrule
\end{tabular}
\end{center}

\subsection{Physical constants}

\begin{center}
\begin{tabular}{@{}lll@{}}
\toprule
Symbol & Value & Defined as \\
\midrule
$c$  & $299\,792\,458$ m/s           & Speed of light \\
$h$  & $6.626\,070\,15 \times 10^{-34}$ J\,s & Planck's constant \\
$k$  & $1.380\,649 \times 10^{-23}$ J/K       & Boltzmann's constant \\
\bottomrule
\end{tabular}
\end{center}

\end{document}
